%!TEX TS-program = xelatex
%!TEX options = -aux-directory=Debug -shell-escape -file-line-error -interaction=nonstopmode -halt-on-error -synctex=1 "%DOC%"
\documentclass{article}
\input{LaTeX-Submodule/template.tex}

% Additional packages & macros
\usepackage{changepage} % Modify page width
\usepackage{multicol} % Use multiple columns
\usepackage[explicit]{titlesec} % Modify section heading styles
\setitemize{leftmargin=*,topsep=1ex,partopsep=0ex,itemsep=-1ex,partopsep=0ex,parsep=1ex}
\setlist{leftmargin=*,topsep=1ex,partopsep=0ex,itemsep=-1ex,partopsep=0ex,parsep=1ex}

\DeclareMathOperator*{\argmin}{arg\,min}
\DeclareMathOperator{\sgn}{sgn}
\usepackage{subcaption}
\usepackage{tabularray}
\usepackage{mdframed}

\titleformat{\section}{\raggedright\normalfont\bfseries}{}{0em}{#1}
\titleformat{\subsection}{\raggedright\normalfont\small\bfseries}{}{0em}{#1}

%% A4 page
\geometry{
	a4paper,
	margin = 10mm
}

%% Hide horizontal rule
\renewcommand{\headrulewidth}{0pt}
\fancyhead{}

%% Hide page numbers
\pagenumbering{gobble}

%% Multi-columns setup
\setlength\columnsep{4pt}

%% Paragraph setup
\setlength\parindent{0pt}
\setlength\parskip{0pt}

% Metadata for README
\newcommand{\unitName}{Machine Learning}
\newcommand{\unitTime}{Semester 1, 2024}
\newcommand{\unitCoordinator}{Dr Simon Denman}
\newcommand{\documentAuthors}{Tarang Janawalkar}

%% Copyright
\usepackage[
    type={CC},
    modifier={by-nc-sa},
    version={4.0},
    imagewidth={5em},
    hyphenation={raggedright}
]{doclicense}

\newenvironment{aside}
  {\begin{mdframed}[style=0,%
      leftline=false,rightline=false,leftmargin=2em,rightmargin=2em,%
          innerleftmargin=0pt,innerrightmargin=0pt,linewidth=0.75pt,%
      skipabove=7pt,skipbelow=7pt]\small}
  {\end{mdframed}}


\begin{document}
% Modify spacing
\titlespacing*\section{0pt}{1ex}{1ex}
\titlespacing*\subsection{0pt}{1ex}{1ex}
%
\setlength\abovecaptionskip{8pt}
\setlength\belowcaptionskip{-15pt}
\setlength\textfloatsep{0pt}
%
\setlength\abovedisplayskip{1pt}
\setlength\belowdisplayskip{1pt}
\begin{multicols}{3}
    \section{ML and Data}
    \textbf{Supervised} Learn map between input and labelled output data:
    predict output for unseen input.

    \textbf{Unsupervised} Knowledge discovery or anomaly detection from input data.

    \textbf{Dataset} is large and diverse. Suffers from \textit{curse of
        dimensionality} due to large \(p\) and small \(n\).
    Sparse feature spaces overfit. Increase in \(p\) necessitates
    complex models, which requires more data.

    \textbf{Training} Training model.\\
    \textbf{Validation} Tuning hyperparameters.\\
    \textbf{Test} Evaluating model.

    If data is limited, use \textit{cross-validation} to dynamically
    split data (80--20) so that each split is used in validation, then
    choose best model.

    \textbf{Normalisation} Scale between \(\interval{0}{1}\).

    \textbf{Standardisation} Mean 0, variance 1.

    \textbf{Encoding} Convert categorical data to numerical.

    \textbf{Mean-Centring} Subtract mean from data.
    \section{Supervised Learning}
    \subsection{Linear Regression}
    \(x^{\left( i \right)}\)-predictor, \(y^{\left( i \right)}\)-response.
    % \begin{align*}
    %     y^{\left( i \right)} = \hat{y}\left( x^{\left( i \right)} \right) + \epsilon^{(i)} \\
    %     \hat{y}\left( x \right) = \sum_{i = 0}^p \hat{\beta}_i x_i = \symbf{\beta}^\top \symbf{x}
    % \end{align*}
    % \(\symbf{\beta}\) found by minimising cost:
    % \begin{align*}
    %     J\left( \symbf{\beta} \right) & = \frac{1}{n} \sum_{i = 1}^n L\left( \symbf{\beta} \right)                                                        \\
    %     L\left( \symbf{\beta} \right) & = \frac{1}{2n} \norm*{\symbf{y} - \symbf{X} \symbf{\beta}} = \frac{1}{2n} \symbf{\epsilon}^\top \symbf{\epsilon}.
    % \end{align*}
    Strength of linear relationship between \(x\) and \(y\) using
    covariance \(s_{xy}\):
    \begin{itemize}
        \item \(s_{xy} > 0\): \(y\) increases as \(x\) increases
        \item \(s_{xy} < 0\): \(y\) decreases as \(x\) increases
        \item \(s_{xy} \approx 0\): no linear relationship
    \end{itemize}
    To generalise, use normalised correlation coefficient \(r_{xy}\) between
    \(-1\) and \(1\).\ \(r\) has no information about gradient and \(r = 0\)
    only implies no \textit{linear} relationship.

    \textbf{Linearity} Relationship is indeed linear.
    Scatter plot (\(\hat{y}\) vs.\ \(y\)).

    \textbf{Multicollinearity} \(x^{\left( i \right)}\) are independent.
    \(\abs*{r_{x_1,x_2}} > 0.7\) suggests high correlation.

    \textbf{Independence} \(y^{\left( i \right)}\) are independent.

    \textbf{Normality} \(\epsilon^{\left( i \right)} \sim \mathcal{N}\left( 0,\: \sigma^2 \right)\).
    QQ-plot and histogram of \(\epsilon\).

    \textbf{Homoscedasticity} \(\epsilon^{\left( i \right)}\) have constant variance.
    Residual plot (\(\epsilon\) vs.\ \(\hat{y}\)).

    \textbf{Exogeneity} No correlation between \(x^{\left( i \right)}\) and \(\epsilon^{\left( i \right)}\).
    Residual plot (\(\epsilon\) vs.\ \(x\)).

    \textbf{R\({}^2\)} Proportion of variance explained by predictors.

    \textbf{Adjusted R\({}^2\)} Adjusts for \#predictors.

    \textbf{F-statistic} Overall significance of model (<0.05).

    \textbf{p-value} Significance of \(\beta_i\) (<0.05).
    Misleading if assumptions are violated.

    \textbf{Regularisation} prevents overfitting by regularising coefficients.
    % \begin{equation*}
    %     J_\lambda\left( \symbf{\beta} \right) = J\left( \symbf{\beta} \right) + \lambda R\left( \symbf{\beta} \right).
    % \end{equation*}
    \textbf{Bias} is error introduced by erroneous assumptions in model.
    High bias \textit{underfits} training data. Increasing complexity reduces bias.

    \textbf{Variance} is error introduced by model's sensitivity to fluctuations in data.
    High variance \textit{overfits} training data. Reducing complexity reduces variance.

    \textbf{LASSO} (L1) \(\beta_i \to 0\): feature selection. Faster during testing.

    \textbf{Ridge} (L2) Reduces large \(\beta_i\), smoother cost function. Faster during
    training.
    \subsection{Classification}
    Predicting classes of categorical data.

    \textbf{SVM (Binary)} separates classes with max margin. Use kernel
    when non-linear. Optimal classifier fails when not
    linearly separable and is sensitive to outliers. Slack variables
    allow misclassification. Small box constraint \(C\) allows more
    misclassifications.

    \textbf{KNN} classifies points using majority class
    of \(k\) nearest neighbours.
    No training phase, sensitive to size of data and outliers. Small
    \(k\) leads to poor performance on noisy data. Large \(k\) leads to
    misclassification when points are close to decision boundary or
    classes imbalanced. Distance metric is problem dependent.
    Distance weighting can be applied to give more weight to closer
    points.

    \textbf{Random forest} combines output of multiple decision trees
    using random subsets of training data. Data is split by class purity
    (Gini inpurity) or information gain. Tree depth should consider
    class imbalance; large tree depths often overfit. More trees provide
    better average predictions.

    \textbf{One-vs-One} pairs class with another class (\(p \left( p + 1 \right)/2\) classifiers).

    \textbf{One-vs-All} pairs class with all other classes (\(p\) classifiers).
    Suffers from class imbalance---use weights inversely proportional to size
    of class.

    \textbf{Evaluation} \(
    \begin{bmatrix}
        \text{TP} & \text{FN} \\
        \text{FP} & \text{TN}
    \end{bmatrix}
    \)
    \begin{itemize}
        \item Recall: minimise false negatives.
        \item Precision: minimise false positives.
        \item Accuracy: overall performance.
        \item F1-score: harmonic mean of precision and recall.
    \end{itemize}
    Consider class imbalance as metrics may be biased toward
    majority class.\ \textbf{Select model} using grid-search to tune
    hyperparameters.
    \section{Neural networks}
    \textbf{Loss functions} decide how errors are penalised during training:
    \begin{itemize}
        \item \textbf{Regression} MSE, MAE
        \item \textbf{Classification} Binary CE (Logistic Loss), Hinge Loss;
              Categorical CE (one-hot), Sparse Categorical CE (integer)
    \end{itemize}
    Output specified as probability or logits (after or before softmax).

    \textbf{Activation} adds non-linearity.

    \textbf{Backpropagation} updates weights to minimise loss using
    gradient descent optimisers.\ \textbf{Epochs} represent time taken
    for one gradient descent step on training set:
    \(\text{updates per epoch} = n/b\).

    Number of epochs consider if validation accuracy continues to
    increase, and if training accuracy is significantly higher than
    validation accuracy (overfitting). \textbf{Small batches} take
    suboptimal paths toward minimum as batches are less representative
    of dataset; may lead to underfitting when classes are imbalanced.
    \textbf{Large batches} may require more epochs for satisfactory
    result and consume more memory.

    \textbf{Ensemble methods} leverage multiple models to improve performance
    but increase computational cost.

    \textbf{Regularisation} performed to prevent overfitting by using dropout
    layers, batch normalisation (layers trained on same scale, improves
    training speed), weight regularisation (L1/2 on weights, biases,
    activation functions), or fine tuning existing models.
    \textbf{Augmentation} increases size of training set. Should not change
    meaning of image, nor be too extreme that it is unlikely.
    \subsection{Convolutional Neural Networks}
    Learn spatial features using kernels. Stack 1/2 convolution, max
    pooling and batch normalisation layers with increasing filter size,
    followed by fully connected layers.\ \textbf{Residual
    networks} introduce skip connections from input layer to output of
    convolution layers (fast training when deep). When number of layers
    (parameters) is large, use bottleneck layers to keep internal
    representation small (\(1 \times 1\) filter at start/end of
    convolution block, or in skip branch: patterns across channels).
    Addresses vanishing gradient problem.

    \textbf{Metric Learning} learns embeddings
    for verification and identification, where
    similar classes are close, and dissimilar classes are far. Can add
    new classes after training. Backbone network is an encoder that is
    trained to learn embeddings.
    \begin{itemize}
        \item Siamese passes pairs of images through network. Uses
              binary CE loss (does not force similar images to be
              close) or contrastive loss.
        \item Triplet passes triplets (anchor, positive, negative) of
              images through network. Uses triplet loss (similar to
              contrastive loss).
    \end{itemize}
    Contrastive and triplet loss separate pairs by a margin.
    Normalise embedding to use a margin of 1.
    Distribution of distances should minimise overlap.\
    \textbf{Embedding size} should be sufficiently
    large and depends on number of classes and complexity of backbone.
    \section{Unsupervised Learning}
    \textbf{K-Means clustering} finds \(k\) clusters that minimise
    distances between data points and their cluster centres without overlap.
    Run multiple times as clusters are sensitive to initialisation of cluster centres
    (pronounced when number of data points/clusters increases).
    Distance metric is problem dependent but restricts algorithm to roughly
    spherical clusters. Scale dimensions to avoid domination.
    \textbf{GMMs} assign a probability that data belongs to
    a cluster. Initialise with \(k\)-means. Computationally expensive.

    \(k\)-means is only used for exploration, as it doesn't measure cluster
    likelihood. GMMs can be used for exploration and detecting abnormalities.
    \(k\)-means is more efficient and scales better than GMMs.

    Determine \(k\) using reconstruction error (average distance to
    assigned cluster; elbow; only for \(k\)-means) or BIC (model
    informativeness; minimum).

    \textit{Under-clustering} true clusters merge into super-cluster.

    \textit{Over-clustering} true cluster divided into sub-clusters.

    \textbf{Dimensionality reduction}: good for sparse data; reduces
    computational cost by removing less important features.

    \textbf{PCA} projects data onto orthogonal components that maximise
    variance. Used to extract features and compress data.
    Choose the top \(q\) features using cumulative sum of explained
    variance (components are sorted by importance).\ \(q = p\) retains
    all information. Must standardise data before applying PCA and
    ensure \(n \gg p\).

    \textbf{LDA} (supervised) finds projection
    that best discriminates between two classes. Returns \(k - 1\) components.
    Sensitive to number of samples per class: each class should have more
    samples than number of features. Large \(n\) or \(p\) leads to poor
    performance, due to overfitting. First apply PCA.\ Unlike PCA, LDA
    cannot reconstruct original data.

    \textbf{t-SNE} (supervised) visualise data in 2D
    using locally linear relationships.

    \textbf{Autoencoders} reconstruct input data (regression/MSE). The
    encoder compresses inputs to a lower-dimensional bottleneck using
    convolution and max pooling with decreasing filter size.

    The decoder then reconstructs input using convolution and
    upsampling with increasing filter size. Consider dimension of input
    for bottleneck.

    \textbf{Multi-Task learning} uses common encoder to train multiple
    networks. Weight loss functions for each task to prevent
    task from dominating. Each task regularises the others. Ensure
    correct loss is used for each task.

    \textbf{Semi-Supervised learning} uses both labelled and unlabelled
    data in training. Loss function does not include unlabelled data.

    \textbf{Variational autoencoders} generate data by learning a
    continuous latent space which we can sample from using a decoder.
    Use KL divergence to constrain distribution to be normal (samples
    distributed across latent space).
    \subsection{Sequences and Attention}
    Sequence data is time-dependent and has variable length.

    \textbf{RNNs} process sequential data using feedback loops, where
    previous time-step information is passed to the network with the
    next time-step. Suffers from vanishing/exploding gradient problem.

    \textbf{LSTM} uses sigmoid and tanh functions to control flow of
    information. Long-term memories are updated using forget gate.
    Short-term memories are updated using input and output gates.
    Placed after embedding layer. May be stacked to improve performance,
    but can be expensive; first LSTM layer should return the sequence to
    the next LSTM but the last LSTM should return a vector of dimension
    appropriate for sequence.
    \begin{itemize}
        \item Sequence-to-one: sequence input and single output: \(X_{t+1} \mid X_t, \ldots, X_0\).
        \item Sequence-to-sequence: sequence input and sequence output: \(X_{t+T}, \ldots, X_{t+1} \mid X_t, \ldots, X_0\).
            Can be used for regression or classification.
    \end{itemize}
    \textbf{Bi-directional} LSTMs process data forward and backwards
    through the network to learn from past and future.

    \textbf{Attention} mechanism allows network to focus on important
    elements of a sequence. Attention weights are learned during training.
    Placed between encoder and decoder. Activation function can be
    softmax or sigmoid, depending on how individual parts of the
    sequence should be weighted. Attention weights can be visualised to
    understand what the network is focusing on.

    \textbf{Transformers} use multi-headed self- attention to measure the
    importance of each element in a sequence in relation to all other
    elements. Multi-headed refers to learning multiple ways to learn
    attention. Outperform LSTMs, but require far more operatios. Can be
    stacked. Outputs are always sequence-to-sequence, but can use dense
    layers or global average pooling to reduce sequence to vector.

    Transformers can be parallelised unlike RNNs and LSTMs which are sequential.
\end{multicols}

\end{document}
